\section{Projecting the clock scores}\label{sec:clock-scores}

We now prove the derivative estimates for the stopped histories.  The
ordinary likelihood coefficients of a path of length $D$ have size
$D^{l/2}$ at order $l$.  We gain in the contribution of a raw coefficient
to a tested expectation by compensating its last ring before taking absolute
values.  For the sharper estimates, we first average an earlier raw block
conditional on its endpoints while later insertions are still raw; those
endpoints are then in the past of the later block in the chosen orientation.
A drift obtained after conditioning on a future endpoint is not thereby
predictable for a smaller filtration.

Throughout this section the interaction matrix is fixed on
\eqref{hist:static-conditions}, and all weights and tested laws are at this
base matrix.  They are common to every $e\in\mathscr E_n$.  Only the ring
likelihood $L_{I,e}(u)$ varies when a clock coefficient is formed.  Write
\[
 \ell_{l,e}(I)=[u^l]L_{I,e}(u),\qquad \ell_{0,e}(I)=1.
\]
Superscripts $\to$ and $\leftarrow$ specify forward and reversed time
orientation.  The reversed object records the reversed marks, including
the old sign as the reversed new sign.
Whenever a bound written $\lesssim_\beta$ absorbs a factor $n^\eta$ from
a path estimate, choose $\eta$ beforehand as a sufficiently small fixed
multiple of $\beta$.  Only a number of such losses depending on $K$ occurs
in this section.

\subsection{Changes of path law and their compensators}

We will use two kinds of oriented probability law.  The first is a tested
path law in its natural forward filtration.  This filtration contains the
initial state and the marked past; future external endpoints are still
unrevealed.  The second is $\mathsf C_I^{x,y}$ on its first two thirds in
the forward orientation, or on its last two thirds in the reversed
orientation, for fixed $x,y$ satisfying the endpoint bound in
Lemma~\ref{hist:conditional-interior}.  Its reference probability is the
ordinary \emph{unconditioned} prefix from its initial endpoint.  The other
endpoint is a fixed parameter, and no additional future path information
is revealed.  An independent history, or a conditional copy sampled using
only already revealed data and independently of the subsequent path, can
be included among the information available at the start.  These will be
called the admissible oriented laws.
When such auxiliary data are adjoined, use the same conditional sampling
kernel under the reference and changed laws, independently of future rings
given the information already present.  This leaves both the reference
intensities and the likelihood ratio unchanged.

A bridge expectation in an identity is taken at its fixed boundary pair.
Every high-probability assertion about its pathwise bounds, however, is
with respect to the tested law of the boundary followed by that conditional
interior.  It is not a uniform conditional assertion for all pairs with
$\bar k_I(x,y)\ge n^{-L}$: the overlap and count events are inherited under
this joint law.  This convention applies to every admissible bridge portion
below.

In either orientation the reference predictable intensity of the mark
$(i,\zeta)$, $\zeta\in\{-1,1\}$, is
\begin{equation}\label{score:reference-rates}
 r_{i,\zeta}(t)=\frac{1+\zeta t_i(X(t-))}{2}.
\end{equation}
Their sum at a site is one.  Under an admissible law $\mathsf Q$, write its
predictable intensities as $r_{i,\zeta}b_{i,\zeta}$.  Such predictable
ratios exist on histories of positive $\mathsf Q$ probability by the
intensity change formula \cite[Theorem~4.1]{Jacod1975}.  For example,
conditional waiting and mark densities under a probability absolutely
continuous with respect to the ordinary path law give them by dividing
the conditional waiting density by the conditional no-ring survival.
Localization covers zeros of that survival.  For a time interval $H$ in
the oriented portion put
\begin{equation}\label{score:hellinger-cost}
 \mathcal J_H=\int_H\sum_{i,\zeta}r_{i,\zeta}(t)
                    (\sqrt{b_{i,\zeta}(t)}-1)^2\,dt,
 \qquad
 \mathcal J_{H,i}=\int_H\sum_\zeta r_{i,\zeta}
                    (\sqrt{b_{i,\zeta}}-1)^2\,dt.
\end{equation}
The cost is additive on disjoint intervals.  The probabilities and hence
the ratios $b_{i,\zeta}$ do not depend on the entry $e$ being differentiated.

\begin{lemma}[Likelihood and intensity bounds]\label{score:hellinger}
Let $\mathcal F_t$ be the natural filtration of a nonexplosive marked path,
augmented by an initial $\sigma$-field, and let $\mathsf Q\ll\mathsf P$ be a change of law up to a
finite stopping time $\tau$, with density $Z$ there and with intensities $rb$ relative to
the reference intensities $r$ conditional on that initial information.
For the cost in \eqref{score:hellinger-cost},
\begin{equation}\label{score:hellinger-inequality}
 \mathsf Q\{\mathcal J_{[0,\tau]}>A,\ \log Z\le B\}
          \le e^{(B-A)/2}.
\end{equation}
Consequently, for every fixed $\eta>0$, with arbitrarily high
inverse-polynomial probability,
\begin{equation}\label{score:action-bounds}
 \begin{cases}
 \mathcal J_{\rm path}\le n^\eta,&\text{under a tested forward law},\\
 \mathcal J_{\rm portion}\le C N(|I|)+n^\eta,
       &\text{under either oriented portion of }\mathsf C_I^{x,y},\quad |I|\ge d_0.
 \end{cases}
\end{equation}
Here $d_0$ is a fixed constant large enough for
Equation~\eqref{hist:prefix-density}.
The second assertion is over the tested boundary law followed by the
conditional interior.  In all these cases the integrated new intensity
at any one site on each specified unit-length block is at most $n^\eta$
with the same probability convention.  These integrated intensities and
$\mathcal J$ have polynomial fixed-order moments, also after averaging
over the boundaries of a conditional interior.
\end{lemma}

\begin{proof}
Include a possible initial density $z_0=d\mathsf Q|_{\mathcal F_0}/
d\mathsf P|_{\mathcal F_0}$.  With bounded integrated rates, the marked
point-process likelihood formula \cite[Theorem~5.1]{Jacod1975} gives
\[
 \sqrt Z\,e^{\mathcal J_{[0,\tau]}/2}
 =\sqrt{z_0}\prod_{s\le\tau\,\mathrm{ring}}\sqrt{b_{i_s,\zeta_s}(s)}
   \exp\left\{-\int_0^\tau\sum_{i,\zeta}r_{i,\zeta}(\sqrt{b_{i,\zeta}}-1)\,dt
       \right\}.
\]
The product allows zero factors; a mark with $b=0$ has no mass under
$\mathsf Q$.  The product-exponential following $\sqrt{z_0}$ is the
stochastic exponential for intensities $r\sqrt b$ relative to $r$.
Its expectation conditional on the initial
data is at most one, as follows by integrating the successive conditional
waiting and mark densities, also when a mark factor is zero; this is the
nonnegative likelihood supermartingale of \cite[Proposition~4.3]{Jacod1975}.
Cauchy--Schwarz
gives $\mathsf E_{\mathsf P}\sqrt{z_0}\le1$.  Thus the right side has
$\mathsf P$ expectation at most one.  For general ratios, stop only when
the integrated reference, changed, and square-root rates reach a finite
level.  Let $Y=\sqrt Z e^{\mathcal J/2}$ on $\{Z>0\}$, with extended
values, and set $Y=0$ on $\{Z=0\}$.  The localized likelihood expressions
converge to $Y$ on $\{Z>0\}$; assigning zero on $\{Z=0\}$ can only decrease
their lower limit.  Fatou's lemma gives $\mathsf E_{\mathsf P}Y\le1$.
No positivity restriction is imposed on $b$ in this localization.
Since $\mathsf Q\{Z=0\}=0$, this expectation bound also makes
$\mathcal J<\infty$ $\mathsf Q$-almost surely.  On the event in
\eqref{score:hellinger-inequality},
$\sqrt Z e^{-\mathcal J/2}\le e^{(B-A)/2}$.  Multiplying this by $Y$
and integrating under $\mathsf P$ proves the inequality.

For a tested forward law, $\log Z\le(M_++M_-)\log n$ by
Definition~\ref{hist:tested-law}.  For a conditional portion,
\eqref{hist:prefix-density} gives $\log Z\le CN(|I|)+L\log n$ outside an
arbitrarily small boundary event.  Since $N(|I|)\ge c n^\beta$, applying
\eqref{score:hellinger-inequality} with these bounds proves
\eqref{score:action-bounds} after increasing $C$.

Let $C_{i,H}$ be the number of rings at site $i$ in a unit block $H$, and
$\Lambda_{i,H}=\int_H\sum_\zeta r_{i,\zeta}b_{i,\zeta}\,dt$.  The
exponential count supermartingale \cite[Proposition~4.3]{Jacod1975} gives
\[
 \mathsf E_{\mathsf Q}\exp\{-C_{i,H}+(1-e^{-1})\Lambda_{i,H}\}\le1,
\]
with the conditional version when the boundary is fixed.  Ordinary unit
counts have every fixed moment bounded; paying a polynomial test density
and using the preserved copy marginals shows $C_{i,H}\le n^{\eta/2}$ with
arbitrarily high inverse-polynomial probability, simultaneously on any
polynomial collection of blocks.  The last exponential bound then gives
$\Lambda_{i,H}\le n^\eta$ with the same precision.  Its tail inequality
also bounds every moment of $\Lambda_{i,H}$ by a constant times a higher
moment of $1+C_{i,H}$.  The conditional assertions here are first averaged
over the boundary law.  Finally,
$(\sqrt b-1)^2\le b+1$ bounds $\mathcal J$ by the sum of the new and ordinary
integrated intensities.  The moment assertions follow, using the polynomial
count moments from Lemma~\ref{hist:crude-ratios}.
\end{proof}

For later use we make the unprojected size and reversal precise.  At a ring
of site $i\in e=\{i,j\}$, the order-$h\ge1$ coefficient of
\eqref{hist:ring-likelihood} is
\begin{equation}\label{score:ring-coefficient}
 \gamma_{h,e}(t,i,\zeta)
       =X_j(t-)^h(\zeta-t_i(X(t-)))(-t_i(X(t-)))^{h-1}.
\end{equation}
It is bounded by $2$, and
$\sum_\zeta r_{i,\zeta}\gamma_{h,e}(t,i,\zeta)=0$.

\begin{lemma}[Raw coefficient size and reversal]\label{score:raw-size}
For every fixed $1\le p<\infty$ and $1\le l\le K$, under ordinary forward
dynamics on an interval of length $d\ge1$,
\begin{equation}\label{score:martingale-moment}
 \left\|\sup_{t\in I}|\ell_{l,e}([\inf I,t])|\right\|_{L^p}
       \le C_{p,K}d^{l/2}.
\end{equation}
For every marked path the exact likelihood reversal is
\begin{equation}\label{score:reversal}
 \frac{L_{I,e}^{\to}(u)}{1+u q_e(X_{\sup I})}
   =\frac{L_{I,e}^{\leftarrow}(u)}{1+u q_e(X_{\inf I})}.
\end{equation}
It follows that, for every fixed $\eta>0$, with arbitrarily high
inverse-polynomial probability under tested laws and their resampled
marginals, simultaneously over entries and specified tree intervals,
\begin{equation}\label{score:raw-high-probability}
 |\ell_{l,e}^{\to}(I)|+|\ell_{l,e}^{\leftarrow}(I)|
       \le n^\eta |I|^{l/2},\qquad 1\le l\le K.
\end{equation}
On a unit-length interval this also holds for partial products, with an
arbitrarily small power allowance, when a bounded fixed-order product of
earlier coefficients is included.
\end{lemma}

\begin{proof}
The order-$l$ prefix coefficient is a martingale starting from zero.  At a
ring its jump is
$\sum_{h=1}^l\gamma_{h,e}\ell_{l-h,e}([\inf I,t))$; the prefix factors
are predictable and every $\gamma_h$ is centered.  The martingale
square-function inequality \cite[Theorem~1.1]{BurkholderDavisGundy1972},
applied at successive ring times, and H\"older's inequality bound its $L^p$
supremum by a constant times
\[
 \sum_{h=1}^l\|C_{e,I}^{1/2}\|_{L^{2p}}
       \left\|\sup_t|\ell_{l-h,e}([\inf I,t])|\right\|_{L^{2p}},
\]
where $C_{e,I}$ is the number of rings at the two sites.  It is Poisson
with mean $2d$, whose fixed moments are $O(d^p)$ for $d\ge1$.  Induction
in $l$, with the same assertion at the higher moment order $2p$, proves
\eqref{score:martingale-moment}.

Detailed balance for marked paths states that the forward path measure
multiplied by its initial stationary mass equals the reversed measure
multiplied by its final stationary mass.  Apply this once at the base
matrix and once using \eqref{hist:stationary-perturbation}, and divide.
This gives \eqref{score:reversal}, including at self-resampling rings.
The Taylor coefficients of the two endpoint factors in that identity are
bounded through order $K$.  Thus the reversed bounds follow from the forward
ones.  Choose $p$ sufficiently large in \eqref{score:martingale-moment} to
pay the polynomial test density, the number of entries and intervals, and
the desired inverse-polynomial precision.  This proves
\eqref{score:raw-high-probability}; resampling preserves the marginal.
On a unit block a fixed-order coefficient is a fixed-degree polynomial in
the two-site count, which gives the last assertion in the same way.
\end{proof}

\subsection{The gain from the last ring}

The compensator identity will always be used before taking absolute values.
If $A_e$ is known before an oriented interval $H$ and all displayed
products are integrable, then the last-ring
expansion and ordinary centering give the exact identity
\begin{equation}\label{score:compensation-identity}
 \begin{split}
 \mathsf E_{\mathsf Q}[A_e\ell_{l,e}(H)]
 =\mathsf E_{\mathsf Q}\sum_{h=1}^l\int_H\sum_{i\in e,\zeta}
     A_e\ell_{l-h,e}([\inf H,t))\gamma_{h,e}(t,i,\zeta)
           r_{i,\zeta}(t)(b_{i,\zeta}(t)-1)\,dt .
 \end{split}
\end{equation}
In reversed time the interval notation is interpreted in that orientation.
The first step is predictable compensation at rates $rb$
\cite[Theorem~2.1]{Jacod1975}; subtracting $r$ uses the
centering of \eqref{score:ring-coefficient}.  All prefix factors are
predictable.  A multiplier depending on an unrevealed future path cannot
be inserted in this identity; a common external test of that future is
already part of the probability $\mathsf Q$ and hence of its intensities.

\begin{lemma}[One compensated insertion]\label{score:one-insertion}
Let $H$ have length $d\ge1$ within an admissible oriented portion of a
checked root.  Suppose $v_i(t)$ is predictable, $|v_i(t)|\le1$, and
does not depend on the index $j$ of the other endpoint or on the new mark
$\zeta$.  Define
\[
 D_{\{i,j\}}(H)=\int_H X_j(t-)v_i(t)
       \sum_\zeta r_{i,\zeta}(b_{i,\zeta}-1)(\zeta-t_i)\,dt
       +\int_H X_i(t-)v_j(t)
       \sum_\zeta r_{j,\zeta}(b_{j,\zeta}-1)(\zeta-t_j)\,dt .
\]
With arbitrarily high inverse-polynomial probability, using the joint
boundary-and-interior convention for a bridge portion,
\begin{equation}\label{score:one-insertion-bounds}
 |D(H)|_2\lesssim_\beta d^{1/4}\sqrt{\mathcal J_H/n},\qquad
 |D(H)|_4\lesssim_\beta d^{5/16}n^{-1/4}\sqrt{\mathcal J_H}.
\end{equation}
For an insertion of order $h>1$, or with bounded predictable weights that
may depend on $e$, the same bounds with $d^{1/2}$ in place of $d^{1/4}$ and
$d^{5/16}$ hold without using overlap estimates.  In particular this
version applies on unit blocks to previously accumulated bounded
coefficients, with their power allowance included.
\end{lemma}

\begin{proof}
For each site write $\alpha_i$ for the signed measure in the first integral
with $X_j$ omitted.  Split it according to $b\le2$ and $b>2$.  On the first
set $|b-1|\le C|\sqrt b-1|$, so Cauchy--Schwarz and $\sum_\zeta r_{i,\zeta}=1$
bound the square integral of its absolute density by $C\mathcal J_{H,i}$.
On the second set $|b-1|\le C(\sqrt b-1)^2$, so its total variation is at
most $C\mathcal J_{H,i}$.  On each unit block this variation is also bounded
by a constant times the new and ordinary integrated site intensities, hence
by $n^\eta$ on the event of Lemma~\ref{score:hellinger}.  If $a_{ik}$ is the sum
of the two variations on the $k$th unit block, these facts imply
\begin{equation}\label{score:block-masses}
 \sum_k a_{ik}^2\le Cn^\eta\mathcal J_{H,i}.
\end{equation}
In particular the large ratios have been treated by total mass, not by a
pointwise square-integrability assertion.

First suppose all times in $H$ have age at least a fixed multiple of $d$
from the \emph{forward} start of its checked root.  Squaring the contribution
of a fixed site $i$ and averaging its other endpoint gives the kernel
$n^{-1}\sum_jX_j(t)X_j(s)=\rho_2(X(t),X(s))$.  By
\eqref{hist:path-overlaps}, on unit block indices this is bounded by
$n^{O(\beta)}(1+|k-k'|)^{-1/2}$.  Adjacent blocks use the bound one.  The
$\ell^1$ sum of this kernel over one index is $O(d^{1/2})$, so
\[
 \frac1n\sum_j\left|\int_H X_j\,d\alpha_i\right|^2
   \le n^{O(\beta)}d^{1/2}\sum_k a_{ik}^2.
\]
Omitting $j=i$ adds at most $n^{-1}(\sum_k a_{ik})^2$, which is absorbed
by the right side since $d\le CT$.  Summing over the contributing site and
dividing by $\binom n2$ gives, using \eqref{score:block-masses},
\[
 |D(H)|_2^2\le \frac{n^{O(\beta)}d^{1/2}}n
                    \sum_i\sum_k a_{ik}^2
       \le n^{O(\beta)}d^{1/2}\mathcal J_H/n.
\]

For the fourth power use the kernel
$n^{-1}\sum_j\prod_{h=1}^4X_j(t_h)=\rho_4(X(t_1),\ldots,X(t_4))$.
The age assumption and \eqref{hist:path-overlaps} bound its unit-block
kernel by a constant times
\[
 n^{O(\beta)}\sum_{1\le h<h'\le4}(1+|k_h-k_{h'}|)^{-3/4}.
\]
For each chosen pair the convolution is at most
$Cd^{1/4}\sum_k a_{ik}^2$; the other two $\ell^1$ sums contribute at most
$d\sum_k a_{ik}^2$.  Thus the fourth-power integral is at most
$n^{O(\beta)}d^{5/4}(\sum_k a_{ik}^2)^2$.  The omitted-diagonal term is
$n^{-1}(\sum a_{ik})^4$ and is again absorbed, since $d^{3/4}/n\le Cn^{-1/2}$.
Now sum over $i$ and divide by $\binom n2$.  Since
$\sum_i\mathcal J_{H,i}^2\le\mathcal J_H^2$, this gives
\[
 |D(H)|_4^4\le\frac{n^{O(\beta)}d^{5/4}}n
                   \sum_i\left(\sum_k a_{ik}^2\right)^2
       \le n^{O(\beta)}d^{5/4}\mathcal J_H^2/n.
\]
The fourth root is the second bound.

For a general $H$, partition it into unit blocks near the forward root
start and pieces whose lengths grow geometrically and are at most a fixed
multiple of their ages.  Apply the preceding estimates on each piece.
For $\gamma=1/4$ or $5/16$, Cauchy--Schwarz gives
$\sum_jd_j^\gamma\sqrt{\mathcal J_{H_j}}
 \le C d^\gamma\sqrt{\mathcal J_H}$, because $\sum_jd_j^{2\gamma}=O(d^{2\gamma})$.
This partition uses the forward age even when the compensation orientation
is reversed; it is a deterministic division of the integrals and changes
no filtration.  Finally, replacing both overlap kernels by one gives
the exponent $1/2$ for either norm.  That proof uses only a common upper
bound on the absolute block masses, so bounded $e$-dependent predictable
weights and higher insertion orders are permitted in this last assertion.
\end{proof}

\subsection{Conditional coefficients and the midpoint identity}

For a node $I=[a,b]$ with fixed boundaries, keep its descendant weights
fixed and define the conditional coefficient series by
\begin{equation}\label{score:R-definition}
 R_e(I;u)=\mathsf C_I^{x,y}\left[
                   \frac{L_{I,e}^{\to}(u)}{1+u q_e(y)}\right],
 \qquad R_{l,e}(I)=[u^l]R_e(I;u).
\end{equation}
In particular $R_0=1$.  Reversal \eqref{score:reversal} and the reversal
of the bridge give the same series from the reversed path with endpoint
$x$.  Write $c_{l,e}^I=\mathsf C_I^{x,y}[\ell_{l,e}^{\to}(I)]$ for the
bridge average of the raw coefficient.  Exactly,
\begin{equation}\label{score:raw-conditional}
 c_l^I=R_l(I)+q_e(y)R_{l-1}(I),\qquad l\ge1.
\end{equation}
For a reversed raw coefficient its end in this formula is the original
start.  These definitions keep all descendant weights fixed.  Their
derivatives are handled separately through \eqref{hist:pre-ratio}.

At the midpoint $v=(a+b)/2$, set $z=X_v$.  With $x,y$ fixed, $z$ has the
base midpoint probability \eqref{hist:midpoint-probability}, and the two
interiors are independent conditional on $z$.  Reversal
\eqref{score:reversal} gives the pathwise factorization
\[
 \frac{L_{I,e}^{\to}(u)}{1+u q_e(y)}
 =(1+u q_e(z))
   \frac{L_{I^-,e}^{\to}(u)}{1+u q_e(z)}
   \frac{L_{I^+,e}^{\leftarrow}(u)}{1+u q_e(z)}.
\]
Conditioning the two interiors on $z$ therefore gives the exact identity
\begin{equation}\label{score:midpoint-series}
 R_e(I;u)=\sum_z\pi_I^{x,y}(z)(1+u q_e(z))
                     R_e(I^-;u)R_e(I^+;u).
\end{equation}
Here the right factor is written in its reversed orientation, and the
midpoint probability is the fixed base one.  The pathwise factorization
shows why the endpoint normalizations leave exactly the single factor
$1+u q_e(z)$.

For $l\ge2$, taking coefficients and grouping the two raw terms gives
\begin{equation}\label{score:midpoint-coefficient}
 \begin{split}
 R_l(I)=\mathsf E_{\pi_I^{x,y}}\bigg[
    c_l^{I^-}+c_l^{I^+,\leftarrow}
    +\sum_{\substack{a+b=l\\a,b\ge1}}R_a(I^-)R_b(I^+)
    +q_e(z)\sum_{\substack{a+b=l-1\\a,b\ge1}}R_a(I^-)R_b(I^+)
                         \bigg].
 \end{split}
\end{equation}
No second or higher midpoint mass coefficient occurs.  This cancellation
is the reason for retaining the endpoint denominator in
\eqref{score:R-definition}.

The two linear terms and the product terms in this recursion will be
estimated differently.  For the linear terms, the conditional interior
identity gives the vector equalities
\begin{equation}\label{score:raw-before-midpoint-norm}
 \begin{aligned}
 \mathsf E_{\pi_I^{x,y}}c_l^{I^-}
   &=\mathsf C_I^{x,y}[\ell_l^{\to}(I^-)],\\
 \mathsf E_{\pi_I^{x,y}}c_l^{I^+,\leftarrow}
   &=\mathsf C_I^{x,y}[\ell_l^{\leftarrow}(I^+)].
 \end{aligned}
\end{equation}
These identities integrate the midpoint back out before compensation.
The right-hand sides use the parent oriented laws with $x,y$ fixed and
$z$ neither conditioned on nor revealed in advance; the natural path
reveals its midpoint state when it reaches that time.  We take vector norms
only after the compensation identity.  The product terms instead use the
conditional child $R$ bounds under the tested boundary law followed by the
conditional interior.  In particular, the estimate of a linear term does
not replace the norm of its midpoint average by an average of raw child norms.

Before estimating this recursion we give the raw estimate that is also
used for a final block of a tested forward path.

For a node $J$ in an oriented subtree, write $J_{\mathrm{early}}$ and
$J_{\mathrm{late}}$ for the children encountered first and second in that
orientation.  They are $(J^-,J^+)$ in forward time and $(J^+,J^-)$ in
reversed time; $J^-,J^+$ continue to denote the chronological children in
the kernel and midpoint identities.  In the raw-representation arguments
below we suppress the arrows on $\ell$ and $c$ when the orientation is fixed.
The state $X_{\mathrm{end}}$ in
$c_j=R_j+q_e(X_{\mathrm{end}})R_{j-1}$ is then the boundary state at the
end of the relevant interval in that orientation.

\begin{lemma}[Representations after compensating the last score]\label{score:raw-projection}
Fix an admissible oriented law $\mathsf Q$, its oriented filtration, and
any permitted auxiliary-data augmentation.  Fix a dyadic subtree $H$ of
length $d$ in that oriented portion and an order $1\le l\le K$.  There is
a finite family of time-integral vectors $V_\alpha^{(l)}(H)$, chosen from
this fixed setup without reference to a subsequent choice of $A$, such that
for every vector $A_e$ measurable at the oriented start of $H$ for which
the displayed products are integrable, coordinatewise,
\begin{equation}\label{score:raw-representation}
 \mathsf E_{\mathsf Q}[A_e\ell_{l,e}(H)]
       =\mathsf E_{\mathsf Q}\left[A_e\sum_\alpha V_{\alpha,e}^{(l)}(H)\right].
\end{equation}
The vectors and $A$ may depend on the same initial information.  With arbitrarily high
inverse-polynomial probability, using the joint boundary-and-interior
convention for a bridge portion, their sum of norms obeys
\begin{equation}\label{score:raw-projection-bounds}
 \begin{split}
 \sum_\alpha|V_\alpha^{(l)}(H)|_2
       &\lesssim_\beta d^{(l-1)/2+1/4}\sqrt{\mathcal J_H/n},\\
 \sum_\alpha|V_\alpha^{(l)}(H)|_4
       &\lesssim_\beta d^{(l-1)/2+5/16}n^{-1/4}\sqrt{\mathcal J_H}.
 \end{split}
\end{equation}
Thus earlier vector factors can be included by H\"older's inequality after
the identity.  This is a representation inside expectation, not a
pointwise assertion about the raw coefficient $\ell_l(H)$.
\end{lemma}

\begin{proof}
For $l=1$, \eqref{score:compensation-identity} is the representation, and
Lemma~\ref{score:one-insertion} gives both bounds.  For $l\ge2$, expand a raw
coefficient by the smallest dyadic node containing all its positive-order
ring insertions.  If all insertions lie in a leaf, keep that leaf term.
Otherwise, at that unique node $J$ the term is
$\ell_{j,e}(J_{\mathrm{early}})\ell_{l-j,e}(J_{\mathrm{late}})$ for some
$1\le j<l$.  This is an exact partition of the ring expansion, including
rings carrying Taylor order
larger than one.

On a leaf, expand according to the last ring and its assigned order as in
\eqref{score:compensation-identity}.  The remaining prefix coefficients
are predictable and are bounded by a power allowance on this unit block,
by Lemma~\ref{score:raw-size}.  The unit-block conclusion of
Lemma~\ref{score:one-insertion} applies.  At a split of length $y$, the earlier
coefficient is known before the later child and is at most
$n^\eta y^{j/2}$ by \eqref{score:raw-high-probability}.  Apply the inductive
representation to the later child, then include this earlier factor by its
supremum norm.  At any one scale the later children are disjoint and number
at most $Cd/y$.  Consequently
\[
 \sum_{J:\,|J|\asymp y}\sqrt{\mathcal J_{J_{\mathrm{late}}}}
      \le C\sqrt{(d/y)\mathcal J_H}.
\]
For the $L^2$ or $L^4$ bound the power before this summation is
$y^{(l-1)/2+\gamma}$, with $\gamma=1/4$ or $5/16$.  Relative to the claimed
right side, this scale thus has the factor
\[
 (y/d)^{(l-2)/2+\gamma}.
\]
It is geometrically summable for $l\ge2$.  The leaf terms cost at most
$C\sqrt{d\mathcal J_H}$ with the respective factor $n^{-1/2}$ or
$n^{-1/4}$, and are also bounded by the claimed expression.  All
compensations used the expectation identity with the earlier coefficient
still predictable.  After absolute values, each representation vector is
bounded by a fixed polynomial in raw clock counts and prefix coefficients,
and in
$\int_H\sum_{i,\zeta}r_{i,\zeta}|b_{i,\zeta}-1|
 \le\int_H\sum_{i,\zeta}r_{i,\zeta}b_{i,\zeta}+n|H|$.
For the sharper representations below, an earlier endpoint projection adds
a conditional coefficient with the uniform coarse bound of
Lemma~\ref{hist:crude-ratios}.  That lemma, the intensity moments in
Lemma~\ref{score:hellinger}, and H\"older's inequality therefore give
polynomial fixed moments for the representation norms, including a fixed
number of copies.  The conditional exceptional-set rule then applies to
each of the polynomially many terms and proves the stated probability
formulation.
\end{proof}

\begin{proposition}[Conditional clock scores and raw representations]\label{prop:clock-scores}
\textit{(i) Conditional coefficients.}
For every deterministic matrix on \eqref{hist:static-conditions}, every
tested path law, and every node $I$ of length $D$ in its checked roots, the
vectors in \eqref{score:R-definition} satisfy, with arbitrarily high
inverse-polynomial probability over their boundary states,
\begin{equation}\label{score:R-basic}
 |R_l(I)|_\infty\lesssim_\beta D^{l/2},\qquad
 |R_l(I)|_2\lesssim_\beta D^{l/2-1},\qquad
 |R_l(I)|_4\lesssim_\beta D^{l/2-1/2},\qquad 1\le l\le K.
\end{equation}
Also,
\begin{equation}\label{score:R-l1}
 |R_3(I)|_1\lesssim_\beta D^{-1/4},\qquad
 |R_l(I)|_1\lesssim_\beta D^{(l-4)/2},\quad 4\le l\le K.
\end{equation}
If $D\ge Tn^{-0.005}$, the $L^4$ bound for $R_1$ has the additional factor
$n^{-0.01}$, the $L^2$ bounds for $R_2,R_3$ have the additional factor
$n^{-0.001}$, and the $L^1$ bound for $R_4$ has the additional factor
$n^{-0.0007}$.  The conclusions hold simultaneously for any polynomial
collection of specified nodes and any fixed number of conditional copies,
with the same exceptional-set convention.

\textit{(ii) Raw representations.}
Fix an admissible oriented law, its oriented filtration, and any permitted
auxiliary-data augmentation, as in Lemma~\ref{score:raw-projection}.  Then
fix a dyadic subtree $H$ of length $d$ in that portion and an order
$1\le l\le K$.  The following representations satisfy
\eqref{score:raw-representation} for every earlier vector $A$ in that
lemma's class, with the representation chosen from this fixed setup before
$A$.  If $l\ge2$ and $d\ge Tn^{-0.02}$, there is a representation for which
\begin{equation}\label{score:raw-improvement}
 \sum_\alpha|V_\alpha^{(l)}(H)|_2
       \lesssim_\beta n^{-0.004}
                   d^{(l-1)/2+1/4}\sqrt{\mathcal J_H/n}.
\end{equation}
For every $1\le l\le K$, there is a representation for which
\begin{equation}\label{score:raw-l1-projection}
 \sum_\alpha|V_\alpha^{(l)}(H)|_1
       \lesssim_\beta \max\{\sqrt d,d^{l/2-5/4}\}\sqrt{\mathcal J_H/n}.
\end{equation}
For $l=4$ and $d\ge Tn^{-0.006}$, the latter estimate has the additional
factor $n^{-0.001}$.  The two families here may differ from each other and
from the family giving the two basic bounds in Lemma~\ref{score:raw-projection}.
These pathwise norm estimates have arbitrarily high inverse-polynomial
probability with that lemma's scope: for a bridge portion the probability
is under the tested boundary law followed by the conditional interior.
They are not uniform conditional assertions at every fixed boundary pair.
\end{proposition}

\begin{proof}
All conditional estimates in this proof are first established under the
tested boundary law followed by its interior sampling.  We then use the
conditional exceptional-set rule following Lemma~\ref{hist:crude-ratios}.
The polynomial coefficient and representation moments proved in
Lemmas~\ref{hist:crude-ratios} and~\ref{score:raw-projection} permit any
required stronger working precision.
This yields the stated high-probability bounds on conditional expectations,
without asserting them for every possible boundary.

\paragraph{Order one and its stronger $L^4$ bound.}
The infinity bound follows at once by applying
\eqref{score:raw-high-probability} inside
\eqref{score:R-definition} and expanding the bounded endpoint denominator.
For $R_1$, split instead at a variable $v$ in the middle third of $I$.
The pathwise reversal identity gives
\[
 R_1(I)=\mathsf C_I^{x,y}\bigl[
       \ell_{1}^{\to}([a,v])+\ell_{1}^{\leftarrow}([v,b])-q_e(X_v)\bigr].
\]
Average this identity uniformly over that middle third.  Compensating the
two ring terms in the parent forward and reversed portions uses bounded
deterministic time weights, common for the other-site index.  By
\eqref{score:action-bounds} and \eqref{score:one-insertion-bounds} their
$L^2$ cost is at most
\[
 n^{O(\beta)}D^{1/4}\sqrt{N(D)/n}
       =n^{O(\beta)}D^{-1/2}.
\]
For the boundary average $Q_e=(3/D)\int_{a+D/3}^{a+2D/3}q_e(X_v)\,dv$,
the exact edge sum
$\sum_{i<j}q_{ij}(X_v)q_{ij}(X_w)
 =\tfrac12(n^2\rho_2(X_v,X_w)^2-n)$ gives
\[
 |Q|_2^2\le\frac{C}{D^2}\int_{a+D/3}^{a+2D/3}\int_{a+D/3}^{a+2D/3}
                       \rho_2(X_v,X_w)^2\,dv\,dw
       \lesssim_\beta D^{-1}.
\]
The last step uses the square of \eqref{hist:path-overlaps}, whose kernel
is $(1+|v-w|)^{-1}$ on this middle third, and absorbs $\log(2+D)$.
This proves the $L^2$ bound for $R_1$.  Interpolation with its infinity
bound gives the ordinary $L^4$ estimate.

We will retain a stronger intermediate $L^4$ improvement.  Its ring terms
can be bounded directly by the second estimate in
\eqref{score:one-insertion-bounds}, giving
\[
 n^{O(\beta)}D^{5/16}n^{-1/4}\sqrt{N(D)}
       =n^{1/4+O(\beta)}D^{-7/16}.
\]
If $D\ge Tn^{-0.03}$ this is at most $n^{-0.02+O(\beta)}$.
The boundary term has infinity norm at most one and $L^2$ norm
$\lesssim_\beta D^{-1/2}$, so its $L^4$ norm is even smaller.  Hence
\begin{equation}\label{score:R1-intermediate}
 |R_1(I)|_4\lesssim_\beta n^{-0.02},\qquad D\ge Tn^{-0.03}.
\end{equation}
For bounded $D$, where the conditional prefix estimate was not invoked,
Lemma~\ref{score:raw-size} gives all the ordinary bounds after a constant change.

\paragraph{Joint induction for the basic bounds and the sharper raw $L^2$ representation.}
We next prove the $L^2$ and $L^4$ bounds by induction on $l$, together with
\eqref{score:raw-improvement}.  Assume \eqref{score:R-basic} in orders below
$l\ge2$.  To prove the raw improvement when $d\ge Tn^{-0.02}$, keep the
basic representation of Lemma~\ref{score:raw-projection} on split scales
$y<dn^{-0.02}$.  The scale factor in that lemma's proof is at most
$(y/d)^{1/4}$, so these scales gain
$n^{-0.005}$.  Their leaf remainder has a still larger gain.

On a larger split, project the earlier raw coefficient \emph{before}
compensating anything in the later child.  In the expectation at this
stage, the remaining raw insertions are outside the earlier child.  By
Lemma~\ref{hist:conditional-interior}, for any such exterior-measurable
integrand $H_e$ the exact identity is
\begin{equation}\label{score:earlier-projection}
 \mathsf E_{\mathsf Q}[\ell_{j,e}(J_{\mathrm{early}})H_e]
       =\mathsf E_{\mathsf Q}
          [c_{j,e}^{J_{\mathrm{early}}}(X_{\partial J_{\mathrm{early}}})H_e].
\end{equation}
For a prefix law this follows by first conditioning on the full exterior
and then taking the prefix marginal.
The new factor is known at the start of $J_{\mathrm{late}}$.  Its $L^4$ norm is at most
$n^{O(\beta)}y^{j/2-1/2}$ by the induction hypothesis and
\eqref{score:raw-conditional}; for $j=1$ the term $q_eR_0$ is simply
bounded by one.  Apply the basic $L^4$ representation to the later
child and pair these factors in $L^2$.  Compared with the basic split
$L^2$ estimate the gain is
\[
 n^{1/4}y^{-7/16}.
\]
Indeed the new power of $y$ is
$j/2-1/2+(l-j-1)/2+5/16=l/2-11/16$, while the old one is $l/2-1/4$.
Here $y\ge dn^{-0.02}\ge Tn^{-0.04}$, so the last ratio is at most
$n^{-0.024+o(1)}$.  Summing these scales and the smaller scales proves
\eqref{score:raw-improvement}, with ample room for the stated exponent.
The identity \eqref{score:earlier-projection} was applied while the later
coefficient was still raw.  It is not applied to a drift from an earlier
compensation, which could retain the earlier interior in its intensities.

Consider now \eqref{score:midpoint-coefficient}.  For each linear raw term,
use \eqref{score:raw-before-midpoint-norm} and apply
Lemma~\ref{score:raw-projection} on the corresponding half in the
\emph{parent} oriented filtration with just $x,y$ fixed.  After the
compensation identity, use the triangle inequality and the pathwise norm
bounds.  The action bound for
that parent portion is $\mathcal J\lesssim_\beta N(D)$, so its $L^2$ cost
is
\[
 n^{O(\beta)}D^{(l-1)/2+1/4}\sqrt{N(D)/n}
       =n^{O(\beta)}D^{l/2-1}.
\]
For the product terms in \eqref{score:midpoint-coefficient}, the induction
hypothesis and H\"older give
$|R_aR_b|_2\le |R_a|_4|R_b|_4\lesssim_\beta D^{(a+b)/2-1}$.
The terms with $q_e$ have the still smaller exponent $(l-1)/2-1$.
Using the conditional exceptional-set rule proves the $L^2$ bound at order
$l$.  Interpolating it with the infinity bound proves its $L^4$ bound.

For clarity we track the low-order margins needed below.  When
$D\ge Tn^{-0.019}$, a half has length at least $Tn^{-0.02}$ for large $n$.
Use \eqref{score:raw-improvement} on the two raw order-two terms.  Its
product term uses \eqref{score:R1-intermediate} on both children.  Thus
\begin{equation}\label{score:R2-intermediate}
 |R_2(I)|_2\lesssim_\beta n^{-0.004},\qquad
 |R_2(I)|_4\lesssim_\beta n^{-0.002}D^{1/2},
             \qquad D\ge Tn^{-0.019},
\end{equation}
where the second bound follows by interpolation.  Similarly, if
$D\ge Tn^{-0.018}$, the raw order-three terms use
\eqref{score:raw-improvement}; their product terms use
\eqref{score:R1-intermediate} and \eqref{score:R2-intermediate} on the
children.  The terms with $q_e$ have an additional factor $D^{-1/2}$
relative to the target bound.  This gives
\begin{equation}\label{score:R3-intermediate}
 |R_3(I)|_2\lesssim_\beta n^{-0.0015}D^{1/2},
                          \qquad D\ge Tn^{-0.018}.
\end{equation}
These deductions use the induction hypothesis at strictly smaller orders,
so the proof of \eqref{score:R-basic} and of these improvements is complete.

\paragraph{The raw $L^1$ representation.}
With the basic conditional bounds established, we prove
\eqref{score:raw-l1-projection} by induction on $l$.  In every split of the
smallest-node decomposition, first apply \eqref{score:earlier-projection}.
In the term $R_j$ of $c_j=R_j+q_eR_{j-1}$, use its $L^2$ bound
$y^{j/2-1}$ and the basic later-child $L^2$ representation.  The power
before the summation at scale $y$ is $y^{l/2-5/4}$.  The term $q_eR_{j-1}$
has an extra $y^{-1/2}$ when $j\ge2$ and is bounded the same way.  Summing
the square roots of the actions at this scale contributes $\sqrt{d/y}$.
Thus the sum over scales is bounded by
\[
 C\sqrt d\sum_{\substack{y\le d\\y\ \mathrm{dyadic}}}
       y^{l/2-7/4}\sqrt{\mathcal J_H/n}
 \le C\max\{\sqrt d,d^{l/2-5/4}\}\sqrt{\mathcal J_H/n},
\]
with a logarithm allowed at an equality of exponents.

When $j=1$, the remaining term in $c_1$ is $q_e$ alone.  It is known before
the later child.  Use \eqref{score:raw-l1-projection} inductively for the
later order $l-1$; order one follows from the $L^2$ estimate in
Lemma~\ref{score:one-insertion}.  Its scale sum is at most
\[
 C\sqrt d\max\{1,d^{l/2-9/4}\}\sqrt{\mathcal J_H/n},
\]
which is no larger than the claimed expression.  Leaf terms cost
$C\sqrt{d\mathcal J_H/n}$ as in Lemma~\ref{score:raw-projection}.  This proves
the ordinary bound at every order.  Throughout, the earlier conditional
coefficient is left inside expectation until the later compensation is
performed, where it is predictable.

At order four the leading scale contribution, divided by its value at
$y=d$, is $(y/d)^{1/4}$.  If $y<dn^{-0.01}$ this gains $n^{-0.0025}$.
If $y\ge dn^{-0.01}$ and $d\ge Tn^{-0.006}$, then
$y\ge Tn^{-0.016}$.  For an earlier order $j=2$ or $3$, its $R_j$ term has
the improved $L^2$ bound \eqref{score:R2-intermediate} or
\eqref{score:R3-intermediate} on that child.  Its $q_eR_{j-1}$ term has
the extra $y^{-1/2}$ just noted.  For $j=1$ with its $R_1$ term, the later
order is three, and \eqref{score:raw-improvement} applies there.  Constant
halves are absorbed by the slack in these thresholds.  Finally the leaf
and $q_e$-alone branches cost only $O(\sqrt d)$ in place of $d^{3/4}$.
Each case gains at least $n^{-0.001+O(\beta)}$.  This proves the additional
factor in \eqref{score:raw-l1-projection}.

\paragraph{Conditional $L^1$ bounds.}
Apply the $L^1$ representation to each raw term in
\eqref{score:midpoint-coefficient}, using the order of conditioning in
\eqref{score:raw-before-midpoint-norm}.  Since
\begin{equation}\label{score:node-action-scale}
 \sqrt{N(D)/n}=n^{\beta/2}D^{-3/4},
\end{equation}
the two raw order-three terms are $\lesssim_\beta D^{-1/4}$, and raw order
$l\ge4$ is $\lesssim_\beta D^{(l-4)/2}$.  For the product terms use the
two $L^2$ estimates in \eqref{score:R-basic}.  At order three they cost at
most $D^{-1/2}$ (and the term with $q_e$ at most $D^{-1}$); at order
$l\ge4$ they cost at most $D^{(l-4)/2}$, with an extra $D^{-1/2}$ in the
terms containing $q_e$.  This proves \eqref{score:R-l1}.

If $D\ge Tn^{-0.005}$, a half satisfies the threshold for the improved
order-four representation.  In its product terms, $R_1R_3$ uses the
improved $L^2$ bound for $R_3$, $R_2R_2$ uses that for $R_2$, and the terms
with $q_e$ already have the extra $D^{-1/2}$.  The claimed factor
$n^{-0.0007}$ follows.  The stated weaker late factors for $R_1,R_2,R_3$
follow immediately from \eqref{score:R1-intermediate}--\eqref{score:R3-intermediate}.
The conditional exceptional-set rule and a union over the specified nodes
and fixed number of copies complete the proof.
\end{proof}

The representations in Lemma~\ref{score:raw-projection} and
Proposition~\ref{prop:clock-scores}(ii) also apply to a last differentiated
root of a tested forward path.  Each is used for its respective norm bound.
In that use \eqref{score:action-bounds} gives
$\sqrt{\mathcal J_H}\le n^\eta$ after adjusting $\eta$.  Earlier root
coefficients can first be conditioned on their endpoints by
\eqref{score:earlier-projection}; these endpoints are then in the past of
the last root.  This is the form used in the replacement argument.
